\documentclass[../../thesis_main.tex]{subfiles}
\ifSubfilesClassLoaded{\externaldocument{../../thesis_main}}{}
\graphicspath{{\subfix{../../figures/}}}

\begin{document}

    \ifSubfilesClassLoaded{\tableofcontents}{}

    \chapter{Materials and Methods}\label{cap2}
    This chapter describes the experimental apparatus used in this work.\CCgen{capitolo 2}{l'apertura non fa la mappa del capitolo}{L'introduzione e' un capoverso solo, e sopra le import restano un capoverso commentato sull'obiettivo a lungo termine e il promemoria ``\%! Penso di metterlo nell'introduzione'': quella decisione va chiusa. Al di la' di questo, il capitolo attraversa cinque cose molto diverse (MOT e vuoto, procedure di misura e calibrazioni, iniezione nella black fiber, breadboard superiore e allineamento, modello numerico) senza dire in apertura quali sono e perche' stanno in quest'ordine. Il capitolo 1 questa mappa ce l'ha. Qui basterebbero tre righe che dicono le cinque parti e cosa serve a cosa, e sarebbe anche il posto naturale per la distinzione fra apparato preesistente e lavoro nuovo di cui alla nota qui accanto.} First, a brief overview of the implementation of the Magneto-Optical Trap (MOT) for rubidium-87 atoms ($^{87}\mathrm{Rb}$) is presented (this setup was already available prior to the present work). The chapter then focuses on the experimental and numerical contributions carried out during this thesis. Section 2.2 describes the calibration of the imaging systems and the experimental procedures used to characterize the cold atomic cloud, including the loading dynamics into the Optical Dipole Trap (ODT). A central focus of this experimental work is the implementation of the optical conveyor belt, which requires the precise coupling and alignment of two counter-propagating laser fields. Section 2.4 details the injection of the 1064~nm trapping light into a polarization-maintaining single-mode fiber (referred to as the ``black fiber''), while Section 2.5 characterizes the power and polarization stability of this beam using Allan-deviation and frequency-domain shot-noise analyses. Section 2.6 discusses the optical layout on the upper breadboard, where the top-down black-fiber beam and the bottom-up beam emerging from the Hollow-Core Photonic Crystal Fiber (HCPCF) are spatially overlapped to form the moving optical lattice. Section 2.7 presents a full multi-level numerical model, developed using QuTiP, to simulate the Electromagnetically Induced Transparency (EIT) cooling dynamics that are required to achieve sub-Doppler temperatures once the atoms are transported inside the hollow core.
    % Although beyond the scope of the current work, it is worth noting that the long-term objective of this experiment is to confine the atoms within the hollow core of the HCPCF, employing a dual beam ODT, to implement an optical conveyor belt.
    %! Penso di metterlo nell'introduzione
    \import{./}{cap2_MOT.tex}
    \clearpage
    \import{./}{cap2_PROCEDURE.tex}
    \clearpage
    \import{./}{cap2_BLACK_setup.tex}
    \clearpage
    \import{./}{cap2_upper.tex}
    \clearpage
    \import{./}{cap2_EIT.tex} 

    \ifSubfilesClassLoaded{\printbibliography}{}
\end{document}